% =====================================================================
% Class-S parent for the K3 x E chiral BKM construction.
% Primary anchors: Gaiotto 2012 (arXiv:0904.2715);
% Chacaltana--Distler 2010 (arXiv:1008.5203);
% Shapere--Tachikawa 2008 (arXiv:0804.1957);
% Beem--Rastelli 2015 and Beem--Lemos--Liendo--Peelaers--
% Rastelli--van Rees 2015.
% =====================================================================

\providecommand{\kBKM}{\kappa_{\mathrm{BKM}}}
\providecommand{\kch}{\kappa_{\mathrm{ch}}}
\providecommand{\kcat}{\kappa_{\mathrm{cat}}}
\providecommand{\kfib}{\kappa_{\mathrm{fiber}}}
\providecommand{\Ntwo}{\mathcal{N}{=}2}
\providecommand{\cS}{\mathcal{S}}
\providecommand{\R}{\mathbb{R}}
\providecommand{\C}{\mathbb{C}}
\providecommand{\Z}{\mathbb{Z}}

\section*{Class-\texorpdfstring{$\mathcal S$}{S} parent for the
\texorpdfstring{$K3\times E$}{K3xE} chiral BKM}
\label{sec:gaiotto-curve-k3e-resolution}

\subsection*{Pants arithmetic for
\texorpdfstring{$T[A_1,\Sigma_{0,24}]$}{T[A1,Sigma0,24]}.}

For the $A_1$ class-$\mathcal S$ theory on a punctured sphere with
maximal regular $\mathfrak{su}(2)$ punctures, the pants decomposition
has $n-2$ trinions $T_2$ and $n-3$ gauged $SU(2)$ tubes.  Each trinion
is the free trifundamental half-hypermultiplet of $SU(2)^3$ and
contributes $(n_v,n_h)=(0,4)$; each tube contributes an $SU(2)$ vector
multiplet, hence $(n_v,n_h)=(3,0)$.  Thus
\begin{equation}
  n_v(\Sigma_{0,n}) = 3(n-3),
  \qquad
  n_h(\Sigma_{0,n}) = 4(n-2).
  \label{eq:sphere-pants-nvnh}
\end{equation}
The Shapere--Tachikawa and Beem--Rastelli relation between four- and
two-dimensional central charges is
\begin{equation}
  c_{4d}=\frac{2n_v+n_h}{12},
  \qquad
  c_{2d}=-12\,c_{4d}.
  \label{eq:class-s-c4d-c2d}
\end{equation}
For $n=24$ this gives
\begin{equation}
  n_v=21\cdot3=63,
  \qquad
  n_h=22\cdot4=88,
  \label{eq:sigma-0-24-nvnh}
\end{equation}
and therefore
\begin{equation}
  c_{4d}\big(T[A_1,\Sigma_{0,24}]\big)
  =\frac{2\cdot63+88}{12}
  =\frac{107}{6},
  \qquad
  c_{2d}=-214.
  \label{eq:sigma-0-24-central-charges}
\end{equation}
Equivalently,
\begin{equation}
  c_{4d}(A_1,\Sigma_{0,n},\mathrm{max\,reg})
  =\frac{5n-13}{6},
  \qquad
  n=24.
  \label{eq:sphere-central-charge}
\end{equation}

\subsection*{Closed genus two gives a different number.}

The same pants count for a closed genus-two $A_1$ curve has two
trinions and three gauged $SU(2)$ tubes:
\begin{equation}
  n_v(\Sigma_{2,0})=3\cdot3=9,
  \qquad
  n_h(\Sigma_{2,0})=2\cdot4=8.
  \label{eq:sigma-2-0-nvnh}
\end{equation}
Hence
\begin{equation}
  c_{4d}\big(T[A_1,\Sigma_{2,0}]\big)
  =\frac{2\cdot9+8}{12}
  =\frac{13}{6},
  \qquad
  c_{2d}=-26.
  \label{eq:sigma-2-0-central-charges}
\end{equation}
This closed curve has \(c_{4d}=13/6\). The Schur VOA with
$c_{4d}=107/6$ comes from \(\Sigma_{0,24}\). The genus-two object belongs to the Siegel
target: the Humbert divisor $H_1\subset\mathcal A_2$, the Jacobian
map to Siegel space $\mathbb H_2$, and the Borcherds product
producing $\Delta_5$.

\begin{proposition}[Class-\texorpdfstring{$\mathcal S$}{S} parent]
\label{prop:gaiotto-curve-k3e}
The class-$\mathcal S$ parent whose Schur construction has
$c_{4d}=107/6$ and $c_{2d}=-214$ is
\[
  T[A_1,\Sigma_{0,24}],
\]
where $\Sigma_{0,24}$ is the sphere with twenty-four maximal regular
$\mathfrak{su}(2)$ punctures.  The closed genus-two curve
$\Sigma_{2,0}$ has $c_{4d}=13/6$ in the same pants model and does not
produce the required central charge.
\end{proposition}

\begin{proof}
Equations~\eqref{eq:sigma-0-24-nvnh} and
\eqref{eq:sigma-0-24-central-charges} compute the punctured-sphere
theory directly from its $22$ trinions and $21$ gauged tubes.  Equation
\eqref{eq:sigma-2-0-central-charges} computes the closed genus-two
theory from its $2$ trinions and $3$ gauged tubes.  The central charges
are distinct:
\[
  \frac{107}{6}\neq\frac{13}{6}.
\]
Thus the curve supplying the Schur central charge is
$\Sigma_{0,24}$.
\end{proof}

\subsection*{Siegel target and Borcherds lift.}

The two curves live in different moduli problems.  The punctured sphere
$\Sigma_{0,24}$ is the Gaiotto curve for the four-dimensional
$A_1$ class-$\mathcal S$ parent.  The genus-two curve appears after
passing to principally polarised Abelian surfaces: its Jacobian maps to
$\mathcal A_2$, and the Humbert divisor $H_1$ records the product
degeneration of elliptic curves.  This is the Siegel-modular side of
the construction. The compactification curve of the six-dimensional
$(2,0)$ theory is \(\Sigma_{0,24}\).

The Schur-to-Borcherds comparison is therefore the composite
\begin{equation}
  \mathcal I_{\mathrm{Schur}}\big(T[A_1,\Sigma_{0,24}]\big)
  \xrightarrow{\;\operatorname{avg}_{M_{24}}\;}
  \phi^{K3}_{0,1}
  \xrightarrow{\;\mathrm{Borch}\;}
  \Delta_5.
  \label{eq:schur-m24-borcherds-composite}
\end{equation}
The first arrow is the Mathieu-covariant averaging to the K3 weak
Jacobi form; the second is the Borcherds lift.  There is no direct
equality between the Schur index and $\Delta_5$.

\subsection*{Numerical invariants.}

The chiral-BKM object $\mathfrak g_{\Delta_5}$ attached to
$K3\times E$ carries the following separated invariants:
\begin{itemize}\setlength\itemsep{0pt}
  \item Gaiotto class-$\mathcal S$ parent
  $T[A_1,\Sigma_{0,24}]$ with
  $(n_v,n_h)=(63,88)$ and
  $(c_{4d},c_{2d})=(107/6,-214)$;
  \item Siegel target $\mathbb H_2/\mathrm{Sp}_4(\Z)$ with
  $\Delta_5$ obtained by the Borcherds lift of $\phi^{K3}_{0,1}$;
  \item $\kBKM(\Delta_5)=c_1(0)/2=5$ by the Gritsenko--Nikulin
  weight computation, with the Igusa-cusp-form calculation serving as
  consistency evidence for the same Borcherds product;
  \item $\kcat(K3\times E)=0$ by K\"unneth multiplicativity on
  $\chi(\mathcal O_X)$.
\end{itemize}

The arithmetic separation is rigid: the punctured sphere supplies the
Schur central charge, the Siegel genus-two variable supplies the
automorphic target, and the Borcherds lift connects them only through
the $M_{24}$-averaged Jacobi form.

\paragraph*{References.}
Gaiotto 2012 \emph{JHEP} 1208:034 (arXiv:0904.2715);
Chacaltana--Distler 2010 \emph{JHEP} 1011:099 (arXiv:1008.5203),
\S5.14, Table~3;
Shapere--Tachikawa 2008 \emph{JHEP} 0809:109 (arXiv:0804.1957), \S3;
Beem--Lemos--Liendo--Peelaers--Rastelli--van~Rees 2015
\emph{Commun.\ Math.\ Phys.} 336:1359 (arXiv:1312.5344);
Beem--Peelaers--Rastelli--van~Rees 2014 (arXiv:1408.6522);
Conway--Sloane 1988 \emph{Sphere Packings, Lattices, and Groups} \S11;
Eguchi--Ooguri--Tachikawa 2010 \emph{Exp.\ Math.} 20;
Gaberdiel--Hohenegger--Volpato 2012 \emph{Commun.\ Number Theory Phys.}
6:1;
Gritsenko--Nikulin 1998 \emph{Amer.\ J.\ Math.} 119;
Lorgat 2020, \S3.1, as a consistency computation for
$\phi^{K3}_{0,1}\mapsto\Delta_5$.

\begin{thebibliography}{9}
\bibitem{ShapereTachikawa2008} A.\ Shapere, Y.\ Tachikawa,
  \emph{Central charges of $\Ntwo$ superconformal field theories in four
  dimensions}, JHEP 0809 (2008) 109, arXiv:0804.1957.
\bibitem{BeemRastelli2015} C.\ Beem, L.\ Rastelli, \emph{Vertex
  operator algebras, Higgs branches, and modular differential
  equations}, JHEP 2018, arXiv:1707.07679.
\end{thebibliography}
